\begin{problem}[第35届全国中学生物理竞赛复赛][穆斯堡尔效应]{95}
八、1958 年穆斯堡尔发现的原子核无反冲共振吸收效应（即穆斯堡尔效应）可用于测量光子频率极微小的变化，穆斯堡尔因此荣获 1961 年诺贝尔物理学奖。类似于原子的能级结构，原子核也具有分立的能级，并能通过吸收或放出光子在能级间跃迁。原子核在吸收和放出光子时会有反冲，部分能量转化为原子核的动能（即反冲能）。此外，原子核的激发态相对于其基态的能量差并不是一个确定值，而是在以 $E_0$ 为中心、宽度为 $2\Gamma$ 的范围内取值的。对于 $^{57}\text{Fe}$ 从第一激发态到基态的跃迁， $E_0 = 2.31 \times 10^{-15}\text{J}$ ， $\Gamma \square 3.2 \times 10^{-13}E_0$ 。已知质量 $m_{\text{Fe}} \square 9.5 \times 10^{-26}\text{kg}$ ，普朗克常量 $h = 6.6 \times 10^{-34}\text{J} \cdot \text{s}$ ，真空中的光速 $c = 3.0 \times 10^{8}\text{m} \cdot \text{s}^{-1}$ 。
\begin{subproblem}
忽略激发态的能级宽度，求反冲能，以及在考虑核反冲和不考虑核反冲的情形下， $^{57}$Fe 从第一激发态跃迁到基态发出的光子的频率之差；
\end{subproblem}
\begin{subproblem}
忽略激发态的能级宽度，求反冲能，以及在考虑核反冲和不考虑核反冲的情形下， $^{57}$Fe 从基态跃迁到激发态吸收的光子的频率之差；
\end{subproblem}
\begin{subproblem}
考虑激发态的能级宽度，处于第一激发态的静止原子核 $^{57}Fe^{*}$ 跃迁到基态时发出的光子能否被另一个静止的基态原子核 $^{57}Fe$ 吸收而跃迁到第一激发态 $^{57}Fe^{*}$ （如发生则称为共振吸收）？并说明理由。
\end{subproblem}
\begin{subproblem}
现将 ${}^{57}Fe$ 原子核置于晶体中，该原子核在跃迁过程中不发生反冲。现有两块这样的晶体，其中一块静止晶体中处于第一激发态的原子核 ${}^{57}Fe^{*}$ 发射光子，另一块以速度 $V$ 运动的晶体中处于基态的原子核 ${}^{57}Fe$ 吸收光子。当速度 $V$ 的大小处于什么范围时，会发生共振吸收？如果由于某种原因，到达吸收晶体处的光子频率发生了微小变化，其相对变化为 $10^{-10}$ ，试设想如何测量这个变化（给出原理和相关计算）？
\end{subproblem}
\end{problem}

